% Magic comments - Informa ao compilador algumas regras de execução, como tipo de compilação ou codificação do texto
% !TeX root = template-abnt-ufop-escola-de-minas.tex
% !TeX encoding = UTF-8
% !BIB TS-program = XeLaTex
% !backend = biber

%% abtex2-modelo-trabalho-academico.tex, v-1.9.2 laurocesar
%% Copyright 2012-2014 by abnTeX2 group at http://abntex2.googlecode.com/
%%
%% This work may be distributed and/or modified under the
%% conditions of the LaTeX Project Public License, either version 1.3 of this license or (at your option) any later version.
%% This work has the LPPL maintenance status `maintained'.

%% INÍCIO DAS CUSTOMIZAÇÕES PARA A UNIVERSIDADE FEDERAL DE OURO PRETO%%
%%
%% 2024.2.21 09h Bruno Nazario Coelho
%% Altera o padrão de TCC para IC
%% Atualização dos dados pré-textuais
%% Adaptação para o Modelo de Relatório de IC da PROPPI / UFOP
%% Inclusão da Ficha de Avaliação na primeira página
%% 
%% 
%% 
%% 2022.3.30 13h15 Danny Tonidandel
%% Altera as definições de capa, alterando  o comando, inserindo figuras próprias da Universidade e alterando a disposição dos elementos na contra-capa
%% Remove itens de abstract em francês
%% Altera disposição dos elementos na folha de aprovação
%% Altera o conteudo dos capítulos para um arquivo único.
%% Personalização do estilo biblatex-abnt para adequação à normaABNT 6023:2018
% Reseta contadores das notas de rodapé em cada capítulo
% Insere comandos para exibir uma caixa colorida de fundo cinza para notas explicativas, à escolha do autor
%%
%% 2017.5.31 21h13 Danny Tonidandel
%% Altera nome de arquivo de logomarca
%% remove resumos em espanhol e italiano
%% Insere exemplos para elaboração de tabelas
%% Insere tabela de cronograma de atividades (para projeto de pesquisa)
%%
%% FIM DAS CUSTOMIZAÇÕES PARA A UNIVERSIDADE FEDERAL DE OURO PRETO

%% This work consists of the files
% template-abnt-ufop-escola-de-minas.tex
% referencias.bib
%%
% --------------------------------------------
% abnTeX2: Modelo de trabalho Academico (tese de doutorado, dissertacao de
% mestrado e trabalhos monograficos em geral) em conformidade com
% ABNT NBR 14724:2011: Informacao e documentacao - Trabalhos academicos -
% Apresentacao
% ------------------------------------------------------------------------

\documentclass[
	% -- opções da classe memoir --
	12pt,				% tamanho da fonte
	% openright,			% capítulos começam em pág ímpar (insere página vazia caso preciso)
	oneside, % twoside,			% para impressão em verso e anverso. Oposto a oneside
	a4paper,			% tamanho do papel.
	% -- opções da classe abntex2 --
	%chapter=TITLE,		% títulos de capítulos convertidos em letras maiúsculas
	%section=TITLE,		% títulos de seções convertidos em letras maiúsculas
	%subsection=TITLE,	% títulos de subseções convertidos em letras maiúsculas
	%subsubsection=TITLE,% títulos de subsubseções convertidos em letras maiúsculas
	% -- opções do pacote babel --
	english,			% idioma adicional para hifenização
	brazil				% o último idioma é o principal do documento
	]{abntex2}

% ---
% Pacotes básicos
% ---
% Inclusão de gráficos
\usepackage{graphicx}
% tabelas
\usepackage[table,xcdraw]{xcolor}
\usepackage{booktabs}
% pasta de figuras
\graphicspath{{figuras/}}
% extensões permitidas
\DeclareGraphicsExtensions{.pdf,.eps,.svg,.png,.jpg,.bmp}

\usepackage{indentfirst} % para incluir indentação no primeiro parágrafo

\usepackage{microtype} 			% para melhorias de justificação
\usepackage{amsmath,amssymb,unicode-math} % escrita matematica
%\usepackage{relsize} % readicionado comando \part não funciona. Ver: https://github.com/abntex/abntex2/issues/109
\usepackage{verbatim}

% para incluir páginas pdf diretamente no documento
\usepackage{pdfpages}
\usepackage{csquotes}

%%%%%%%%%%%%%%%%%%%%%%%%%%%%%%%
% Pacotes de citações: retire o comentário do estilo que deseja usar
%%%%%%%%%%%%%%%%%%%%%%%%%%%%%%%%
% Usando o pacote biblatex-abnt e o biber para a compilacao das fontes

% Opção 1: notas de referência, com citações reduzidas (idem, ibidem opcit...)
%\usepackage[backend=biber,
%style=abnt-ibid,
%sccite,
%scbib,
%justify,
%extrayear,
%citecount,
%noslsn,
%backref,
%repeatfields]{biblatex}

% Opção 2: notas explicativas no sistema autor-data
\usepackage[backend=biber,
% configuracoes do estilo abnt
style=abnt,
sccite, % sobrenomes em caixa alta
ittitles, % Titulos em italico
citecount, % contar o número de citações
scbib, % biliografia em caixa alta
justify,
noslsn,
repeatfields,
sorting=nty, % ordem alfabetica
]{biblatex}

% ARQUIVO COM AS REFERÊNCIAS BIBLIOGRAFICAS
\addbibresource{referencias.bib}

% ---
% Personalização do estilo biblatex-abnt - por Danny A. V Tonidandel
% Adequa as urls de acordo com normas 6023:2018
\DeclareFieldFormat{url}{\bibstring{urlfrom}\addcolon\addspace \url{#1}}%
\DeclareFieldFormat{urldate}{\bibstring{urlseen}\addcolon\addspace #1}%

% Reseta contadores das notas de rodapé em cada capítulo
\makeatletter
\@addtoreset{footnote}{chapter}
\makeatother

% Pacotes adicionais, usados apenas no âmbito do Modelo Canônico do abnteX2 - podem ser removidos
% ---
% Pacotes adicionais, usados no anexo do modelo de folha de identificação
% ---
% Tabelas
\usepackage{supertabular}
\usepackage{multicol}
\usepackage{multirow}
\usepackage{lipsum}	% para geração de dummy text
% ---


% ---
% Informações de dados para CAPA e FOLHA DE ROSTO
% ---
\titulo{Modelo \LaTeX~ para Projetos de IC - Engenharia de Controle \& Automação - DECAT - UFOP}
\autor{Seu Nome Completo}
\local{Ouro Preto}
\data{\the\year} % imprime o ano corrente no campo data
\orientador{Prof. Bruno Nazário Coelho, PhD.}
\coorientador{Prof. Theophrastus von Hohenheim, PhD.}
\instituicao{Universidade Federal de Ouro Preto}
\tipotrabalho{Relatório Final de Atividades da Iniciação Científica}
\preambulo{Relatório Final, referente ao período (especificar os meses relacionados ao período do relatório), apresentado à Universidade Federal de Ouro Preto, como parte das exigências do programa de iniciação cientifica / (Nome da agência financiadora da bolsa quando houver).}
% ---


% ---
% Configurações de aparência do PDF final

% alterando o aspecto da cor azul
\definecolor{blue}{RGB}{41,5,195}

% informações do PDF
\makeatletter
\hypersetup{
     	%pagebackref=true,
		pdftitle={\@title},
		pdfauthor={\@author},
    	pdfsubject={\imprimirpreambulo},
	    pdfcreator={LaTeX with abnTeX2},
		pdfkeywords={ufop}{latex}{abntex}{decat}{ic},
		colorlinks=true,       		% false: boxed links; true: colored links
    	linkcolor=blue,          	% color of internal links
    	citecolor=blue,        		% color of links to bibliography
    	filecolor=magenta,      		% color of file links
		urlcolor=blue,
		%bookmarksdepth=4
}
\makeatother
% ---

% ---
% Espaçamentos entre linhas e parágrafos
% ---

% O tamanho do parágrafo é dado por:
\setlength{\parindent}{1.3cm}

% Controle do espaçamento entre um parágrafo e outro:
\setlength{\parskip}{0.2cm}  % tente também \onelineskip

% ---
% compila o indice
% ---
\makeindex
% ---

% ----
% Início do documento
% ----
\begin{document}

% Retira espaço extra obsoleto entre as frases.
\frenchspacing

% ----------------------------------------------------------
% ELEMENTOS PRÉ-TEXTUAIS
% ----------------------------------------------------------
% \pretextual

% ---
% Inserir Ficha de Avaliação final de bolsista / voluntário e do Projeto
% ---

% Isto é um exemplo de Ficha de Avaliação final de bolsista, elemento obrigatório da PROPPI / UFOP. Você pode utilizar este modelo até a aprovação do trabalho. Após isso, substitua todo o conteúdo deste arquivo pela ficha gerada pelo seu orientador, fazendo upload e substituindo o arquivo do comando abaixo:
%
\includepdf{fichaavaliacao_final.pdf}

% ---
% Capa
% ---
%\imprimircapa
\begin{capa}
	\thispagestyle{empty}
		\centering
	\begin{center}
		\begin{minipage}{1\linewidth}
			\centering
			\begin{tabular}{ccc}
				\begin{tabular}{c}
					\\
					\includegraphics[width=1.2cm]{figuras/logo-ufop.jpg}
				\end{tabular}
				&
				\begin{tabular}{c}
					\\
					{\large \imprimirinstituicao} \\
					{\large Escola de Minas} \\
					{\large Engenharia de Controle e Automação}
				\end{tabular}
				&
				\begin{tabular}{c}
					\\
					\includegraphics[width=3.0cm]{figuras/logo-EM.jpg}
				\end{tabular}
			\end{tabular}
		\end{minipage}

\centering
\vspace*{3cm}
{\ABNTEXchapterfont\large\imprimirautor}
\vspace*{\fill}

{\ABNTEXchapterfont\bfseries \large \imprimirtitulo}
\vspace*{\fill}

{\large\imprimirtipotrabalho}
\vspace*{\fill}

{\large\imprimirlocal}, {\large\imprimirdata}
%\par
\vspace*{1cm}
	\end{center}
\end{capa}

% ---

% ---
% Folha de rosto
% (o * indica que haverá a ficha bibliográfica)
% ---
\imprimirfolhaderosto*
% ---


% ---
% Inserir a ficha bibliográfica
% ---

% Isto é um exemplo de Ficha Catalográfica, ou ``Dados internacionais de
% catalogação-na-publicação''.
% Porém, caso o relatório seja registrado/publicado, poderá trocar pelo PDF
% com a ficha catalográfica. Quando estiver com o documento, salve-o como PDF no diretório do seu projeto e substitua todo o conteúdo de implementação deste arquivo pelo comando abaixo que está comentado
% (nao se esqueça de comentar o antigo ambiente de ficha catalográfica):
%
% \begin{fichacatalografica}
%     \includepdf{fig_ficha_catalografica.pdf}
% \end{fichacatalografica}
%
%% Ou, você poderá também ler os dados da ficha e adicionar no próprio código
% como as palavras-chave, CDU e dimensões do trabalho:
\begin{fichacatalografica}
	\vspace*{\fill}					% Posição vertical
	\hrule							% Linha horizontal
	\begin{center}					% Minipage Centralizado
	\begin{minipage}[c]{12.5cm}		% Largura

	\imprimirautor

	\hspace{0.5cm} \imprimirtitulo  / \imprimirautor. --
	\imprimirlocal, \the\year-

	  \hspace{0.5cm}  p. : (color.).\\
	\hspace{0.5cm} \imprimirorientadorRotulo~\imprimirorientador\\

	\hspace{0.5cm}
	\parbox[t]{\textwidth}{\imprimirtipotrabalho~--~\imprimirinstituicao,
	\imprimirdata.}\\

	\hspace{0.5cm}
		1. Palavra-chave1.
		2. Palavra-chave2.
		I. Orientador.
		II. Universidade Federal de Ouro Preto.
		III. Escola de Minas.
		IV. Título\\

	% \hspace{8.75cm} CDU 02:141:005.7\\

	\end{minipage}
	\end{center}
	\hrule
\end{fichacatalografica}
% ---


%% ---
%% Dedicatória
%% ---
%\begin{dedicatoria}
%   \vspace*{\fill}
%   \flushright
%   \noindent
%   \textit{Dedico esse trabalho a} \vspace*{\fill}
%\end{dedicatoria}
%% ---

% ---
% Agradecimentos
% ---
\begin{agradecimentos}
\noindent Os agradecimentos vem aqui... não se esqueça de mim :) 

\end{agradecimentos}
% ---

% ---
% Epígrafe
% ---
\begin{epigrafe}
    \vspace*{\fill}
	\begin{flushright}
		\textit{``A possibilidade de realizarmos um sonho é o que torna a vida interessante.'' \\ O Alquimista. (Paulo Coelho)}
	\end{flushright}
\end{epigrafe}
% ---

% ---
% RESUMOS
% ---

% resumo em português
\setlength{\absparsep}{18pt} % ajusta o espaçamento dos parágrafos do resumo
\begin{resumo}
 \noindent O resumo deve ressaltar o objetivo, o método, os resultados e as conclusões do documento. A ordem e a extensão
 destes itens dependem do tipo de resumo (informativo ou indicativo) e do tratamento que cada item recebe no documento original. O resumo deve ser precedido da referência do documento, com exceção do resumo inserido no
 próprio documento. (\ldots) As palavras-chave devem figurar logo abaixo do resumo, antecedidas da expressão Palavras-chave:, separadas entre si por
 ponto e finalizadas também por ponto.

 \textbf{Palavras-chaves}: latex. abntex. editoração de texto.
\end{resumo}

% resumo em inglês
\begin{resumo}[Abstract]
 \begin{otherlanguage*}{english}

\noindent This is the english abstract.

   \vspace{\onelineskip}

   \noindent
   \textbf{Key-words}: latex. abntex. text editoration.
 \end{otherlanguage*}
\end{resumo}

% resumo em francês
%\begin{resumo}[Résumé]
% \begin{otherlanguage*}{french}
%    Il s'agit d'un résumé en français.
%
%   \textbf{Mots-clés}: latex. abntex. publication de textes.
% \end{otherlanguage*}
%\end{resumo}
%
%% resumo em espanhol
%\begin{resumo}[Resumen]
% \begin{otherlanguage*}{spanish}
%   Este es el resumen en español.
%
%   \textbf{Palabras clave}: latex. abntex. publicación de textos.
% \end{otherlanguage*}
%\end{resumo}
%% ---
% ---
% inserir lista de ilustrações
% ---
\pdfbookmark[0]{\listfigurename}{lof}
\listoffigures*
\cleardoublepage
% ---

% ---
% inserir lista de tabelas
% ---
\pdfbookmark[0]{\listtablename}{lot}
\listoftables*
\cleardoublepage
% ---

% ---
% inserir lista de abreviaturas e siglas
% ---
\begin{siglas}
  \item[ABNT] Associação Brasileira de Normas Técnicas
  \item[abnTeX] ABsurdas Normas para TeX
\end{siglas}
% ---

% ---
% inserir lista de símbolos
% ---
\begin{simbolos}
  \item[$ \alpha $] Letra grega alpha
  \item[$ \Delta $] Delta
  \item[$ \infty $] Infinito
  \item[$ \Theta $] Letra grega Theta  
  \item[$ \Sigma $] Somatório / Letra grega Sigma  
  \item[$ \in $] Pertence
  \item[$ \Omega $] Letra grega omega
\end{simbolos}
% ---

% ---
% inserir o sumario
% ---
\pdfbookmark[0]{\contentsname}{toc}
\tableofcontents*
\cleardoublepage
% ---
% Comando para resetar contadores das notas de rodapé
\makeatletter
\@addtoreset{footnote}{chapter}
\makeatother



% ----------------------------------------------------------
% ELEMENTOS TEXTUAIS
% ----------------------------------------------------------
\textual
% ---
% Modelo de capitulo com a introducao, objetivos e estrutura do texto
% ---
% ----------------------------------------------------------
\chapter[Introdução]{Introdução}
%\addcontentsline{toc}{chapter}{Introdução}
% ----------------------------------------------------------
Este documento e seu código-fonte são exemplos de referência de uso da classe \textsf{abntex2} e do pacote \textsf{biblatex-abnt}. O documento exemplifica uma realização possível entre as opções existentes na norma ABNT NBR 10520:2018 \emph{Citações em documentos -- Apresentação} e da norma ABNT NBR 6023:2018 \emph{Referências -- Elaboração}, cientes de que existe uma distância entre as ``normas'' e a interpretação das normas. Assim, antes de tudo, converse com seu orientador ou representantes do programa de pós-graduação de sua universidade, mostre uma cópia do documento PDF gerado por este arquivo e certifique-se de que não terá problas futuros com relação à aceitação ou não do modelo.

\section{Justificativas e Relevância}

A expressão ``Modelo Canônico'' é utilizada para indicar que \abnTeX\ não é modelo específico de nenhuma universidade ou instituição, mas que implementa tão somente os requisitos das normas da ABNT.

\section{Objetivos}

O principal objetivo deste documento é detalhar o modelo para o relatório final de Iniciação Científica, facilitando a vida dos meus orientandos e, consequentemente, a minha.

Descrever o objetivo principal do trabalho em uma frase curta e objetiva.

Descrever os objetivos secundários do trabalho, em parágrafo ou em estrutura de tópicos. %comando itemize

\begin{itemize}
	\item item;
	\item item;
	\item item;
	\item item;
\end{itemize}

%\section{Organização e estrutura}

%Detalharem um breve parágrafo, como é a estrutura do seu trabalho

\section{Cronograma}

Esta seção deve constar somente no projeto de pesquisa. Não deve aparecer na versão final do texto. Um exemplo de cronograma das atividades é proposto na tabela~\ref{tab:cronograma}. Você pode elaborar também tabelas online ou a partir de qualquer planilha eletrônica, inclusive em outros estilos, gerando o código em \LaTeX. Após isso, basta copiar e colar o código aqui. Um exemplo de site é o ``Tables Generator'': \url{http://www.tablesgenerator.com/}.

% Cronograma
\begin{table}[htpb] 
    %esta sigla htpb ajuda no posicionamento da tabela ao longo do texto. %Existem outras opções, dê uma pesquisada caso necessário
	\centering
	\caption{Cronograma de atividades.}
	\label{tab:cronograma}
	\resizebox{\textwidth}{!}{%
		\begin{tabular}{@{}|l|llllllllllll|@{}}
			\toprule
			&
			\multicolumn{12}{c|}{\cellcolor[HTML]{FFFFFF}Meses} \\ \midrule
			\cellcolor[HTML]{FFFFFF}{\color[HTML]{343434} Atividades (Etapas)} &
			\multicolumn{1}{l|}{01} &
			\multicolumn{1}{l|}{02} &
			\multicolumn{1}{l|}{03} &
			\multicolumn{1}{l|}{04} &
			\multicolumn{1}{l|}{05} &
			\multicolumn{1}{l|}{06} &
			\multicolumn{1}{l|}{07} &
			\multicolumn{1}{l|}{08} &
			\multicolumn{1}{l|}{09} &
			\multicolumn{1}{l|}{10} &
			\multicolumn{1}{l|}{11} &
			12 \\ \midrule
			\cellcolor[HTML]{FFFFFF}{\color[HTML]{343434} 1. Estudo da teoria} &
			\multicolumn{1}{l|}{X} &
			\multicolumn{1}{l|}{X} &
			\multicolumn{1}{l|}{X} &
			\multicolumn{1}{l|}{X} &
			\multicolumn{1}{l|}{X} &
			\multicolumn{1}{l|}{} &
			\multicolumn{1}{l|}{} &
			\multicolumn{1}{l|}{} &
			\multicolumn{1}{l|}{} &
			\multicolumn{1}{l|}{} &
			\multicolumn{1}{l|}{} &
			\\ \midrule
			\cellcolor[HTML]{FFFFFF}{\color[HTML]{343434} 2. Atualização bibliográfica} &
			\multicolumn{1}{l|}{X} &
			\multicolumn{1}{l|}{X} &
			\multicolumn{1}{l|}{X} &
			\multicolumn{1}{l|}{X} &
			\multicolumn{1}{l|}{X} &
			\multicolumn{1}{l|}{X} &
			\multicolumn{1}{l|}{X} &
			\multicolumn{1}{l|}{} &
			\multicolumn{1}{l|}{} &
			\multicolumn{1}{l|}{} &
			\multicolumn{1}{l|}{} &
			\\ \midrule
			\cellcolor[HTML]{FFFFFF}{\color[HTML]{343434} 3. Seleção de Material} &
			\multicolumn{1}{l|}{X} &
			\multicolumn{1}{l|}{X} &
			\multicolumn{1}{l|}{X} &
			\multicolumn{1}{l|}{X} &
			\multicolumn{1}{l|}{X} &
			\multicolumn{1}{l|}{X} &
			\multicolumn{1}{l|}{X} &
			\multicolumn{1}{l|}{X} &
			\multicolumn{1}{l|}{X} &
			\multicolumn{1}{l|}{} &
			\multicolumn{1}{l|}{} &
			\\ \midrule
			\cellcolor[HTML]{FFFFFF}{\color[HTML]{343434} 4. Desenvolvimento do trabalho} &
			\multicolumn{1}{l|}{} &
			\multicolumn{1}{l|}{} &
			\multicolumn{1}{l|}{} &
			\multicolumn{1}{l|}{} &
			\multicolumn{1}{l|}{X} &
			\multicolumn{1}{l|}{X} &
			\multicolumn{1}{l|}{X} &
			\multicolumn{1}{l|}{X} &
			\multicolumn{1}{l|}{X} &
			\multicolumn{1}{l|}{X} &
			\multicolumn{1}{l|}{X} &
			X \\ \midrule
			\cellcolor[HTML]{FFFFFF}{\color[HTML]{343434} 5. Elaboração de Artigo} &
			\multicolumn{1}{l|}{} &
			\multicolumn{1}{l|}{} &
			\multicolumn{1}{l|}{} &
			\multicolumn{1}{l|}{} &
			\multicolumn{1}{l|}{} &
			\multicolumn{1}{l|}{} &
			\multicolumn{1}{l|}{} &
			\multicolumn{1}{l|}{} &
			\multicolumn{1}{l|}{X} &
			\multicolumn{1}{l|}{X} &
			\multicolumn{1}{l|}{X} &
			X \\ \midrule
			\cellcolor[HTML]{FFFFFF}{\color[HTML]{343434} 6. Entregas finais} &
			\multicolumn{1}{l|}{} &
			\multicolumn{1}{l|}{} &
			\multicolumn{1}{l|}{} &
			\multicolumn{1}{l|}{} &
			\multicolumn{1}{l|}{} &
			\multicolumn{1}{l|}{} &
			\multicolumn{1}{l|}{} &
			\multicolumn{1}{l|}{} &
			\multicolumn{1}{l|}{} &
			\multicolumn{1}{l|}{} &
			\multicolumn{1}{l|}{} &
			X \\ \bottomrule
		\end{tabular}%
	}
\end{table}

% ----------------------------------------------------------
% Capitulo de revisão de literatura
% ----------------------------------------------------------
\chapter{Referencial teórico}
% ---

Para citações, se usar o comando ``cite'', será gerada a entrada \cite{descartes-carta-mersene} ou, se quiser apenas o autor, com o comando ``citeauthor'', será gerada a entrada\citeauthor[][12]{descartes-carta-mersene}. Ainda, se quiser imprimir apenas a data de referida obra, basta utilizar o comando ``citedate'', que gera a entrada \citedate[55]{descartes-carta-mersene}.

Deve ser seguido o mesmo padrão ao longo do texto.

De acordo com \footcite[Descartes foi um grande filósofo francês do século XVII. Veja em][p.~10]{descartes-carta-mersene},  etc... Comecemos então com um exemplo de citação, como esta aqui, feita em notas explicativas,\footcite[Esta é uma nota explicativa. Cf. e.g.,][\S 12]{boyle1772} conforme a norma NBR 10520:2018, existem dois tipos de citações em notas de rodapé:

\begin{citacao}
	\begin{itemize}
		\item[$3.6$] \textbf{notas de rodapé:} Indicações, observações ou aditamentos ao texto feitos pelo autor, tradutor ou editor, podendo	também aparecer na margem esquerda ou direita da mancha gráfica.
		\item[$3.7$] \textbf{notas explicativas:} Notas usadas para comentários, esclarecimentos ou explanações, que não possam ser incluídos no texto.
	\end{itemize}
\end{citacao}

%\caixa{O autor não é obrigado a seguir as notas explicativas com o comando ``footcite'', que gera a essa nota explicativa.\footcite{boyle1772} Pode utilizar o estilo autor-data igualmente, mas deve seguir o mesmo padrão ao longo do texto.}


\section{Separar por áreas de estudo}
% ---

Si lapidem multiplicare volueris, audi. Partem eius accipe, multiplica eam cum duabus partibus hydrargyri animati tui, coque sicut in principio fecisti
matrases post eam in farinam mollem et levem de eodem clibano et calore. Multo minus tempore secundus terminus rotae philosophi et pulvis decies plus virtutis acquirit quam primogenito suo. Conare dare unum nent
vel magis, si vis, et thesaurum inaestimabile compleveris optimum quod est in toto mundo, ultra quod aspirare non potes magis. Valetudinem et opes habebis si ea uteris ut decet. habebis thesaurus qui dat omnem felicitatem mundanam et quod ego pauper rusticus Ego Pontis- sanus, ter feci et feci in domo mea in platea dos Escrivains, prope capellam Sancti-Jacques-de-la-Boucherie; quod ego, Nicolau Flamel, concedo tibi pro amore, quem tibi dedico, et tuo honorem Dei, gloriam et laudem Patris et Filii et Spiritus Sancti Sancta, sancta Trinitas, cui precor ut ex hoc tempore; illuminet et perducat in viam veritatis, lucem et semitam salutem.

% ---
% Capitulo com exemplos de comandos inseridos de arquivo externo
% ---
\chapter{Metodologia e Desenvolvimento}
%
Um bom exemplo de fórmula:

\begin{equation}
u(t) = K_p e(t) + K_i\int_{0}^{t} e(\tau)d\tau + K_d \frac{de(t)}{dt}
\label{eq:01}
\end{equation}

Quem fez a disciplina de Teoria de Controle comigo deve se lembrar bem da equação \eqref{eq:01}, que está na página \pageref{eq:01}, onde temos que...

Os projetos de Iniciação Científica sob minha orientação normalmente são desenvolvidos no Campus ou na Escola de Minas do centro (ver figura \ref{fig:escolademinas}, que está na página \pageref{fig:escolademinas}).

Inserindo uma figura. A figura \ref{fig:AlchLab} ilustra os primórdios de um laboratório de pesquisa.

\begin{figure}[htbp]
	\centering
	\includegraphics[scale=1.5]{figuras/labAlchemy.jpg} % teste os valores
	\caption[Legenda reduzida - aparece no sumario]{Interior de um laboratório de Alquimia. Legenda completa - aqui você pode colocar uma explicação melhor, sem que ela apareça no sumário do seu trabalho. Fonte: \cite[]{alchlab2024}.}
	\label{fig:AlchLab}
\end{figure}

Agora vem uma citação. Com Nicholau Flamel, surgem os primeiros estudos sobre...:\footcite{flamel1414}

Eis uma citação bem simples:
\begin{citacao}
	(...) Ego Nicholaus Flamel, clericus Parisiensis, anno 1414, et regni nostri beatissimi principis nostri Karoli sexti, quem Deus benedixit, et post obitus mei fidelis comitis Perrenelle, recordantis sui; Ex phantasia et satisfactione abstuli me scribere in gratiam tuam, carissime nepotem omne dominium secreti proiectionis pulveris seu Dye Philosophus, qui Deo placuit cum parvo servo dispensare; quod ut facias, si feceris ut indicavero. (...)
\end{citacao}

(...) sequitur; ergo orationibus philosophorum sollertia et intellegentia de secreto, sed non scripta ad litteram, quia etiam si secundum naturam intelligi possint, non erunt Utilis. Noli oblivisci ergo Deum orare ut dimittam te intellectus rationis, veritatis et naturae, ut hoc videas liber, in quo arcanum verbo et pagina scriptum est paginam, sicut feci et elaboravi cum carissima matertera tua Perrenelle, quae Tam vehementer memini. Magisterium igitur posui in librum istum, ut ne obliviscaris magni boni quod Deus tibi et pro quod favet tibi. Hoc vt non deficias in eius memoria, to psallite et psallite. et nihil aptius tam bonum eventum celebrare quam carmina magnifica.

\begin{figure}[htbp]
	\centering
	\includegraphics[width=15cm]{figuras/EMcomIA.png} % teste os valores
	\caption[Legenda reduzida - aparece no sumário]{Escola de Minas editada com ferramenta de Inteligência Artificial. Legenda completa - Aqui você pode colocar uma explicação melhor, sem que ela apareça no sumário do seu trabalho.  
    \\ Fonte: Próprio autor - gerada por IA (Befunky.com) a partir de \cite[]{guiadaestrada2024}.}
	\label{fig:escolademinas}
\end{figure}


% ---
% capitulo de Resultados
\chapter{Resultados}
% ---

\section{Dados, dados, dados}
% ---

Et ideo manu mea scripsi hunc librum, et erat ibi ad Sancti-Jacobi ecclesiam destinata, cum in dicta parrochia, post- quam inveni librum Abrahae iudaei, pecuniam vendere nolui et eam diligentissime custodivi ut sic dicta secreta scriberet in eo alchimiae in litteris et characteribus imaginationis meae, quam tibi do clavem. Cura igitur hoc secretum servare et numquam obliviscere ut taceas, memento mei cum sim in sindone,
memorans, nunc tale tibi documentum paravi, ut tecum fac me magnum magistrum alchimiae philosophicae, quia contuli meae delectatio, desiderium, consolatio et phantasia tale secretum tibi concedunt. Ita facio quod ego feci et adhuc facio aetate provecta et decrepita, et omnia in honore et magisterio alchimiae; per naturam.(...)

% ----------------------------------------------------------
% Finaliza a parte no bookmark do PDF
% para que se inicie o bookmark na raiz
% e adiciona espaço de parte no Sumário
% ----------------------------------------------------------
%\phantompart
% ---
% Insere arquivo de Considerações Finais ou Conclusões
% ---
\chapter{Conclusões} % esse capítulo pode ser numerado ou não \chapter*{...
%\addcontentsline{toc}{chapter}{Conclusão}

Um \emph{grand finale}: Si multum aurum facere vis, carissime nepos, quod nunquam est quid ex incommodo et noxa expediat? centum unciae argenti vivi in ​​magno caldario ferreo super ignem validum.

Cum fumantes, unciam pulveris parati coccinea quartae haustus, cum cera adligat eam sicut a parva pila et iacit in dictum condimentum fumantes. fumus cito peribit. Ignem excita et mox transfigurabis; partim in massa, et partim in puluerem auri flaui, quem in fusorium. In massa vel massam proicies et ex omnibus hoc extrahe hydrargyrum circa 99, 170 uncias auri puri, qualitatis insuperabilis, qua uteris ut tibi videbitur. 

Hic es, carissime nepos, multo ditior omnibus regibus, plus enim habebis
ut ipsi et quid possint per totum regnum suum disponere. Sed non producunt aurum nisi paulatim, prudenter, sine revelatione nichil alicui et sine aliis semper confidere. Declaravi vobis thesaurum omnium thesaurorum hujus mundi, quos habui fecique meis manibus tam carus ac bene dilecto comiti Perrenelle. Hoc igitur thesauro utere in nomine Domini Pater et Filius et Spiritus Sanctus, ut sine sollicitudine vivas; dives in hoc mundo, et palmam gloriae consequaris in regno Dei, quod volo.

% ----------------------------------------------------------
% ELEMENTOS PÓS-TEXTUAIS
% ----------------------------------------------------------
\postextual
% ----------------------------------------------------------

%% ----------------------------------------------------------
%% Referências bibliográficas
%% ----------------------------------------------------------
% toca nome de bibliografia para ``Referências''
\printbibliography[title=Referências]

% ----------------------------------------------------------
% Glossário
% ----------------------------------------------------------
%
% Consulte o manual da classe abntex2 para orientações sobre o glossário.
%
%\glossary

% ----------------------------------------------------------
% Apêndices
% ----------------------------------------------------------
%(Lembre-se: Apendices são de autoria do próprio autor do texto.
% Anexos são elementos de autorias de outros, que o autor do texto julga interessante apresentar)
% ---
% Inicia os apêndices:
% ---
\begin{apendicesenv}

% Imprime uma página indicando o início dos apêndices
\partapendices
% ---
% Insere arquivo com os apêndices A e B
% --------
\chapter{texto adicional do mesmo aluno da pesquisa}

Ut te certiorem facias quomodo hac medicina utaris ad sanitatem de corpore tuo et de memoria, moneo te ut omisso vase cum purpureum est, temperandum est, id est, grana quaedam solve vinum album aut sublimatum, donec vinum coloratur, aureum jus iustum quia torqueo. Noli timere infirmis ministrare 12 vel XV guttae in vino vel sorbitione vel alio potu. Quasi per miraculum;
Mox curabitur. Noli semper de illo gloriari, quia homines sunt malus et invidus, quando non potest idem facere quod alii alii. Tu ergo, ut tuam quotidianam obtineas valetudinem, accipe 9 guttas ejus, soluta et aurea quater anni, id est, die 22° Martii, 22 Iunii, 22 Septembris et 22 Decembris, in liquore quodcunque mavis. Utere, ut te doceam, nunquam laboras
aegritudinibus et vita beata frueris, plena salute et divitiis; tu dominus eris omnis naturae, plures lapides habebis aurum et argentum prae principibus et regibus.

\end{apendicesenv}
% ---

% ----------------------------------------------------------
% Anexos
% ----------------------------------------------------------
%(Lembre-se: Apendices são de autoria do próprio autor do texto.
% Anexos são elementos de autorias de outros, que o autor do texto julga interessante apresentar)
% ---
% Inicia os anexos
% ---
\begin{anexosenv}

% Imprime uma página indicando o início dos anexos
\partanexos

% ---
% Insere arquivo com os anexos 1, 2 e 3
\chapter{elementos de autorias de outras pessoas}
% ---
Unum facere potes, scilicet si vis flores et fructus in frigido hyeme solvit sex grana predicte pollinis rubei ad exitum ipsius vas in X sextarios aquae roris calidae et aqua talis arboris vel floris quicquid vis, intra domum ponens vel cooperiens eos paleas vel faenum. Videbis cito subito et admirabile herbam et incrementum, quod multum mireris. 

% ---
\end{anexosenv}

%---------------------------------------------------------------------
% INDICE REMISSIVO
%---------------------------------------------------------------------
%\phantompart
\printindex
%---------------------------------------------------------------------

\end{document}
